% cloud-vs-on-device.tex — deciding where computation (and ML inference) runs: latency, cost,
% privacy, offline, battery, model/data size, freshness, security of logic, coverage,
% observability; on-device ML (Foundation Models, Core ML / Core AI, Vision; Gemini Nano via
% AICore / ML Kit GenAI on Android) vs Private Cloud Compute vs your own server LLM; hybrid
% patterns (on-device first + server fallback, router/cascade, split pre-processing, server
% trains / device infers, precompute + cache, speculative draft); model delivery; syncing and
% versioning results. A decision tree and a decision table.
% Not repeated (neighbours): the Foundation Models API itself; Core ML formats / compute units /
% compression, Vision, Speech (coreml-vision-on-device).
% Sources (checked 2026-09-25): security.apple.com/blog/private-cloud-compute (five core
% requirements; device wraps the request key only to nodes whose attested measurements are in
% the public transparency log; third-party OHTTP relay hides the source IP);
% developer.android.com/ai/gemini-nano (AICore system service; ML Kit GenAI APIs: Prompt,
% Summarization, Proofreading, Rewriting, Image Description, Speech Recognition; AICore
% distributes and updates the model); developer.apple.com/documentation/backgroundassets (iOS 16+;
% managed asset packs via AssetPackManager, optional Apple hosting; essential before launch,
% prefetch, on demand);
%   machinelearning.apple.com/research/apple-foundation-models-2025-updates (on-device model
%     ~3B parameters, 2 bits per weight via QAT) — checked 2026-10-06
%   developer.apple.com/videos/play/wwdc2026/241 (PrivateCloudComputeLanguageModel needs an
%     entitlement; no cloud API cost under 2M first-time downloads; per-user daily limit, higher
%     with iCloud+; 32,000-token context; contextSize + tokenCount(for:) shipped in iOS 26.4)
%     — checked 2026-10-06
% @labels: area=architecture kind=architecture platform=general topic=ai,system-design,architecture discussion=77
% @tags: on-device-ml, cloud-inference, private-cloud-compute, foundation-models, core-ml, gemini-nano, hybrid-inference, model-delivery, background-assets, privacy, latency, trust-boundary
\documentclass[8pt]{extarticle}
\usepackage{printup-sheet}
\usepackage{array}

\lstdefinelanguage{SwiftSheet}{
  morekeywords={protocol,class,final,struct,enum,func,var,let,init,case,public,static,
    if,else,return,guard,self,Self,nil,try,await,async,throws,throw,private,catch,do,
    true,false,some,any,Sendable,Error,String,Int,Void},
  sensitive=true, morecomment=[l]{//}, morestring=[b]",
  literate={->}{{\hbox{-}\hbox{>}}}2 {==}{{\hbox{=}\hbox{=}}}2 {??}{{\hbox{?}\hbox{?}}}2
           {!=}{{\hbox{!}\hbox{=}}}2 {&&}{{\hbox{\&}\hbox{\&}}}2 {<=}{{\hbox{<}\hbox{=}}}2}
\lstset{basicstyle=\ttfamily\scriptsize, aboveskip=2pt, belowskip=2pt}

\newcommand\ct[1]{\texttt{#1}}
\newcommand\hd[1]{\par\vspace{3pt}\noindent{\bfseries\color{sheetBlue}#1}\par\vspace{1pt}}
\newcommand\yes{\textcolor{sheetGreen!60!black}{\ensuremath{\checkmark}}}
\newcommand\no{\textcolor{sheetRed}{×}}

\tikzset{
  q/.style={box, font=\scriptsize, draw=sheetGrey, fill=black!3, text width=27mm, align=center,
            inner sep=2pt, minimum height=11mm},
  srv/.style={box, font=\scriptsize\bfseries, draw=sheetBlue, fill=sheetBlue!10, text width=27mm,
              align=center, inner sep=2pt},
  dev/.style={box, font=\scriptsize\bfseries, draw=sheetGreen!60!black, fill=sheetGreen!12,
              text width=27mm, align=center, inner sep=2pt},
  hyb/.style={box, font=\scriptsize\bfseries, draw=sheetOrange, fill=sheetOrange!12,
              text width=27mm, align=center, inner sep=2pt},
  lbl/.style={font=\tiny, text=black!75, inner sep=1pt, align=center},
}

\begin{document}

\sheettitle{Cloud vs on-device — where should it run?}{architecture · folio}

\oneliner{Place each computation where its \textbf{binding constraint} is cheapest: \textbf{trust}
(anything authoritative runs on a server — the client is attacker-controlled), \textbf{knowledge
and size} (world knowledge, big models, live data $\to$ cloud), \textbf{privacy, offline and
latency} (personal data, per-frame work, no network $\to$ device). Most real features are
\textbf{hybrid}: device first, server (or Private Cloud Compute) as the fallback, results cached
and \textbf{versioned}.}

\vspace{3pt}
\noindent\begin{tikzpicture}[sheet]
  \node[font=\bfseries\small, text=sheetBlue, anchor=west] at (-1.5,0.9) {Ask in this order — the first ``yes'' decides};
  \node[q] (q1) at (0,0) {\textbf{1} Must the result be \textbf{trusted}? (money, entitlements, fraud, shared state)};
  \node[q] (q2) at (3.35,0) {\textbf{2} Needs \textbf{world knowledge}, live data or a model too big for a phone?};
  \node[q] (q3) at (6.7,0) {\textbf{3} Is the input \textbf{personal} (health, messages, photos, location)?};
  \node[q] (q4) at (10.05,0) {\textbf{4} Must it work \textbf{offline} or per frame / per keystroke?};
  \node[hyb] (h) at (13.45,0) {\textbf{HYBRID}\\\mdseries\tiny device first · server fallback;\\cost + coverage decide};
  \node[srv] (a1) at (0,-1.45) {SERVER\\\mdseries\tiny authoritative; client only displays};
  \node[srv] (a2) at (3.35,-1.45) {SERVER LLM / API\\\mdseries\tiny or PCC for private reasoning};
  \node[dev] (a3) at (6.7,-1.45) {DEVICE\\\mdseries\tiny PCC if it doesn't fit; never your logs};
  \node[dev] (a4) at (10.05,-1.45) {DEVICE\\\mdseries\tiny Core ML / Vision / system LLM};
  \foreach \a/\b in {q1/q2,q2/q3,q3/q4,q4/h} \draw[flow] (\a) -- node[lbl, above]{no} (\b);
  \foreach \a/\b in {q1/a1,q2/a2,q3/a3,q4/a4} \draw[hot] (\a) -- node[lbl, right]{yes} (\b);
  \node[lbl, anchor=north west, align=left, text=sheetRed] at (11.9,-1.0) {a ``yes'' at 1 \textbf{and} 3 (e.g. fraud score on\\private data): server decides, device\\pre-processes and redacts first};
\end{tikzpicture}

\vspace{2pt}
{\footnotesize\setlength\tabcolsep{3pt}
\noindent\begin{tabular}{@{}>{\bfseries\raggedright\arraybackslash}p{19mm}>{\raggedright\arraybackslash}p{47mm}>{\raggedright\arraybackslash}p{43mm}>{\raggedright\arraybackslash}p{50mm}@{}}
\toprule
factor & \textbf{on-device} & \textbf{Private Cloud Compute} (Apple) & \textbf{your server / cloud LLM} \\ \midrule
latency & no network; first model load is slow — load early, reuse & network round trip & round trip + queue; radio wake-up on cellular \\
cost per call & free to you; paid in the user's battery + heat & Apple's servers; quota + entitlement (iOS 27 API) & you pay per request / GPU hour, scales with users \\
privacy & data never leaves the phone & leaves, but stateless, no privileged access, verifiable & leaves: consent, retention, region, subprocessors \\
offline & \yes & \no & \no \\
model / quality & small (system LLM $\approx$3B, 2-bit; your Core ML model) & much larger; 32K-token context, reasoning & any size, frontier models, tools, RAG over your data \\
freshness & updates with the OS / app / asset pack & Apple-managed & instant; live data; world knowledge \\
coverage & AI-capable devices only; memory + thermal limits & Apple Intelligence devices only\unverified & every device, Android, web — one behaviour \\
logic security & \textbf{untrusted}: code + weights extractable, results forgeable & not your server — no custom logic & \textbf{authoritative}; secrets stay here \\
observability & no server logs; on-device evals, opt-in telemetry & none for you & full logs, A/B, evals — mind PII \\
\bottomrule
\end{tabular}}

\begin{multicols}{2}
\raggedright\small\setstretch{1.0}

\hd{The options}
\begin{itemize}
  \item \textbf{Apple on-device}: Foundation Models (system LLM, offline, free; check
        \ct{availability}); Core ML / Core AI for \emph{your} model; Vision, Natural Language,
        Speech for common tasks — try these first.
  \item \textbf{Android on-device}: Gemini Nano runs in the \textbf{AICore} system service; apps
        use \textbf{ML Kit GenAI} APIs (Prompt, Summarization, Proofreading, Rewriting, Image
        Description, Speech Recognition). AICore ships and updates the model — no model in your APK.
  \item \textbf{Private Cloud Compute}: Apple's server extension of on-device AI. Five
        requirements — \textbf{stateless} computation, \textbf{enforceable} guarantees,
        \textbf{no privileged runtime access}, \textbf{non-targetability}, \textbf{verifiable
        transparency}. The device encrypts only to nodes whose \textbf{attested} software is in a
        public \textbf{transparency log}; an \textbf{OHTTP relay} run by a third party hides the IP.
  \item \textbf{Your server}: any model, tools and data; you own cost, keys, logs and legal basis.
        API keys never ship in the app — the app calls \emph{your} backend, which calls the model.
\end{itemize}

\hd{Hybrid patterns}
\begin{itemize}
  \item \textbf{Device first, server fallback}: unavailable device, context too long, low
        confidence, guardrail refusal $\to$ escalate (with consent).
  \item \textbf{Router / cascade}: a cheap on-device classifier decides which requests are
        ``hard'' and need the big model.
  \item \textbf{Split}: device pre-processes — OCR, speech-to-text, embeddings, \textbf{PII
        redaction}, downscaling — and sends less, safer data.
  \item \textbf{Server trains, device infers}: ship or download the model; personalise on device
        (Core ML \ct{MLUpdateTask}); only the model moves, not the data.
  \item \textbf{Precompute + cache}: shared results (recommendations, translations of shared
        content) computed once on the server, cached on devices.
  \item \textbf{Speculative}: show the on-device draft now, replace it if the server's answer
        arrives and differs — label which one the user sees.
\end{itemize}

\hd{Traps}
\begin{itemize}
  \trap{``On-device is free'' — battery, heat, memory (jetsam), app size, device coverage.}
  \trap{Validating a purchase, score or discount on the client — anything on the device can be
        patched; the server decides.}
  \trap{A vendor LLM key in the app ``to save a backend'' — extracted in minutes; proxy it.}
  \trap{Server fallback that silently uploads personal data — fallback needs consent + redaction.}
  \trap{No unavailable path: the system model is missing on older or ineligible devices.}
\end{itemize}

\columnbreak

\hd{Example — device first, server fallback}
\begin{lstlisting}[language=SwiftSheet]
protocol Summarizer: Sendable {
  func summarize(_ text: String) async throws -> String }
struct HybridSummarizer: Summarizer {
  let local: any Summarizer          // Foundation Models session
  let remote: any Summarizer         // YOUR backend, not a vendor key
  let consented: @Sendable () -> Bool
  func summarize(_ text: String) async throws -> String {
    if SystemLanguageModel.default.isAvailable,
       text.count < 6_000 {          // rough fit; 26.4+: tokenCount(for:)
      do { return try await local.summarize(text) }
      catch { /* context / guardrail / refusal: escalate */ }
    }
    guard consented() else { throw SummaryError.needsConsent }
    return try await remote.summarize(redactPII(text))
  }
}
\end{lstlisting}

\hd{Model delivery and results}
\begin{itemize}
  \item \textbf{Bundle} the model (app size, every user pays) \emph{or} \textbf{download} it:
        Background Assets managed asset packs (essential before launch, prefetch, or on demand;
        optional Apple hosting), then \ct{MLModel.compileModel(at:)} for a Core ML file.
  \item Store \textbf{model id + version} beside every output; cache results keyed by
        \ct{hash(input) + modelVersion}; a new model invalidates, never silently mixes.
  \item Results are \emph{data}: sync them through the normal sync path (private $\to$ the
        user's devices only). Compute \textbf{shared} results once on the server, \textbf{private}
        ones per device.
  \item Same feature, different engines $\to$ different answers on iOS, Android, old phones:
        design the UI for variable quality, evaluate each path separately.
\end{itemize}

\hd{Remember}
\textbf{Trust $\to$ server · knowledge $\to$ cloud · private, offline, per-frame $\to$ device ·
everything else: device first, cloud fallback, version the results.}


\end{multicols}

\noindent{\footnotesize\color{sheetGrey}\textit{Related:} app-hardening-privacy · sync-engines (syncing results) ·
app-size-and-energy · distributed-system-design}

\end{document}
